%!TEX root = ../../thesis.tex

\section{Where the Approach Helps}\label{sec:discussion-advantages}

% TODO: Interpretation, not a re-listing of \cref{chap:results}. Every claim here must
% point at a specific result, and claims the results do not support must not appear.
%
% Candidate arguments, each to be kept only if the evidence carries it:
%   - deep, sequentially dependent constraints: where a coverage-guided climber stalls
%     because no single mutation improves coverage, a learned value function that
%     carries credit across steps does not (\cref{sec:results-input-actions}),
%   - sample efficiency where an execution is expensive: under emulation, spending
%     compute per decision to spend fewer executions is the right trade, and that is
%     the regime this domain sits in (\cref{sec:rl-sample-efficiency}),
%   - amortized planning: the operator sweep buys many executions per decision, so
%     the planner's cost is paid down rather than paid per input
%     (\cref{sec:results-corpus-actions}),
%   - protocol-state navigation without a hand-written state machine
%     (\cref{sec:target-open62541}).
%
% State the scope of each claim: which target, which baseline, which currency. An
% advantage demonstrated only where the reward was designed
% (\cref{sec:targets-microbenchmarks}) is a claim about credit assignment, not about
% fuzzing --- keep that distinction visible here rather than only in
% \cref{sec:discussion-threats}.
